\documentclass[11pt]{article}
\usepackage[T1]{fontenc}
\usepackage[utf8]{inputenc}
\usepackage{lmodern,microtype,amsmath,amssymb,amsthm,mathtools}
\usepackage[margin=1in]{geometry}
\usepackage{booktabs,array,longtable,enumitem}
\usepackage[hidelinks]{hyperref}
\hypersetup{pdftitle={Sharp Poisson asymptotics for the longest arithmetic progression in a random set},
pdfauthor={ChatGPT (OpenAI)},pdfsubject={Research manuscript for the ProveIt project},
pdfkeywords={arithmetic progressions, Poisson approximation, asymptotics, entropy, lattice effects}}
\usepackage{fancyhdr}
\pagestyle{fancy}\fancyhf{}
\fancyhead[L]{\small Sharp Poisson asymptotics for random progressions}
\fancyhead[R]{\small Research manuscript}
\fancyfoot[C]{\thepage}
\setlength{\headheight}{14pt}
\setlength{\emergencystretch}{2em}
\newtheorem{theorem}{Theorem}[section]
\newtheorem{lemma}[theorem]{Lemma}
\newtheorem{proposition}[theorem]{Proposition}
\newtheorem{corollary}[theorem]{Corollary}
\theoremstyle{definition}\newtheorem{definition}[theorem]{Definition}
\theoremstyle{remark}\newtheorem{remark}[theorem]{Remark}
\newcommand{\PP}{\mathbb P}
\newcommand{\EE}{\mathbb E}
\newcommand{\Var}{\operatorname{Var}}
\newcommand{\Cov}{\operatorname{Cov}}
\newcommand{\ind}{\mathbf 1}
\newcommand{\AP}{\mathcal A}
\newcommand{\cH}{\mathcal H}
\newcommand{\la}{\lambda}
\newcommand{\eps}{\varepsilon}
\newcommand{\dd}{\,\mathrm d}
\DeclareMathOperator{\Pois}{Poisson}
\title{Sharp Poisson asymptotics for the longest\protect\\arithmetic progression in a random set\\[5pt]
\large An entropy coefficient, a second-order correction,\protect\\and the lattice oscillation of the optimal error}
\author{ChatGPT (OpenAI)\\\small Prepared for Vladimir Reshetnikov's ProveIt research program}
\date{October 7, 2026}
\begin{document}
\maketitle
\begin{abstract}
Let $U_n$ be the longest arithmetic progression in the Bernoulli-$p$ random
subset of $[n]$, with $p\in(0,1)$ fixed. We study the accuracy of the
Poisson approximation to $\PP(U_n<r)$, rather than only its limiting law.
Writing $M_r(n)$ for the number of $r$-term progressions, the exact
mean of progression heads is
$\la_{n,r}=M_r(n)p^r-M_{r+1}(n)p^{r+1}$.
We give a self-contained argument for the uniform expansion
\[
 \PP(U_n<r)=e^{-\la_{n,r}}
 \left(1+c_p\la_{n,r}^2\frac{r^2}{n}\right)
 +O_p\!\left(\frac{(\log n)^{3/2}}{n}\right),
 \qquad 2\le r\le n,
\]
where $c_p=\frac{1-p}{p}(\frac53-\frac{\pi^2}{9})$.
In the critical window the remainder is $O_{p,a,b}(r^{3/2}/n)$ and
there is a corresponding logarithmic expansion. Consequently the
uncorrected uniform error is exactly of order $(\log n)^2/n$.
We further determine its lattice-periodic leading coefficient and its
explicit normalized liminf and limsup. The coefficient arises from the
squared binary-entropy profile of the incidence degrees of long
progressions. A local avoidance expansion for nonmonotone cylinder events,
with a connected-triple remainder, supplies the probabilistic step.
The same results hold for the smooth Poisson parameter used by Zhao and
Zhang. All asserted asymptotic results below have written proofs; finite
exact-arithmetic regressions are supplied separately.
\end{abstract}
\noindent\textbf{Status and scope.}
This is an unrefereed research manuscript, not a Lean-verified result.
The sharp rate, explicit correction, and lattice-error formulas are the
contributions proposed for independent review. The literature search does
not certify priority. The argument does not assume the correctness of any
new manuscript in either linked repository, and it does not improve a
deterministic Szemer\'edi or van der Waerden bound.
\clearpage
\tableofcontents
\newpage

\section{Problem, results, and relation to existing work}
\subsection{The random sequence}
Fix $p\in(0,1)$ and put $q=1-p$ and $\alpha=-\log p>0$.
Let $\xi_1,\ldots,\xi_n$ be independent Bernoulli-$p$ variables.
The random set is $A_n=\{x\in[n]:\xi_x=1\}$. An arithmetic progression
always has positive integer common difference. Let $U_n$ be its largest
possible length in $A_n$, with $U_n=0$ for the empty set and $U_n=1$
when there is no two-term progression but the set is nonempty.
For $2\le r\le n$, write
\begin{align}
 D_r&=\left\lfloor\frac{n-1}{r-1}\right\rfloor,\qquad
 M_r(n)=\sum_{d=1}^{D_r}\bigl(n-(r-1)d\bigr)
       =nD_r-\frac{(r-1)D_r(D_r+1)}2,\label{eq:M}\\
 \la_{n,r}&=p^rM_r(n)-p^{r+1}M_{r+1}(n),\label{eq:lambda}\\
 \bar\la_{n,r}&=\frac{n^2p^r(p+qr)}{2r(r-1)}.\label{eq:barlambda}
\end{align}
The definition of $M_r$ also applies when $r>n$, giving zero.
All logarithms are natural. The notation $O_p$ here denotes a deterministic
implicit constant depending on $p$, not stochastic boundedness. Such constants may
blow up as $p\downarrow0$ or $p\uparrow1$.

The key positive constants are
\begin{equation}
 H=\int_0^1\bigl(-u\log u-(1-u)\log(1-u)\bigr)^2\dd u
   =\frac56-\frac{\pi^2}{18},\qquad
 c_p=\frac{2qH}{p}=\frac qp\left(\frac53-\frac{\pi^2}{9}\right).
 \label{eq:constants}
\end{equation}
In particular, $H\approx0.285021977717258$ and
$c_{1/2}\approx0.570043955434516$.

\subsection{Main statements}
\begin{theorem}[Uniform second-order approximation]\label{thm:uniform}
For every fixed $p\in(0,1)$ and either choice
$\eta_{n,r}=\la_{n,r}$ or $\eta_{n,r}=\bar\la_{n,r}$,
\begin{equation}
 \sup_{2\le r\le n}\left|
 \PP(U_n<r)-e^{-\eta_{n,r}}
 \left(1+c_p\eta_{n,r}^2\frac{r^2}{n}\right)\right|
 =O_p\!\left(\frac{(\log n)^{3/2}}{n}\right).
 \label{eq:uniformcorrected}
\end{equation}
\end{theorem}
The polynomial correction in this statement is intentional. Extending the
exponential $\exp(-\eta+c_p\eta^2r^2/n)$ to very small $r$ would be
incorrect: far outside the critical window its positive quadratic
exponent can become enormous.

\begin{theorem}[Critical-window logarithmic expansion]\label{thm:critical}
Fix $0<a<b<\infty$. Uniformly over integers $r$ satisfying
$a\le\la_{n,r}\le b$,
\begin{align}
 r&=\frac{2\log n-\log\log n+O_{p,a,b}(1)}{\alpha},\label{eq:criticalr}\\
 \log\PP(U_n<r)
 &=-\la_{n,r}+c_p\la_{n,r}^2\frac{r^2}{n}
       +O_{p,a,b}\!\left(\frac{r^{3/2}}n\right),\label{eq:criticallog}\\
 \PP(U_n<r)-e^{-\la_{n,r}}
 &=c_p e^{-\la_{n,r}}\la_{n,r}^2\frac{r^2}{n}
       +O_{p,a,b}\!\left(\frac{r^{3/2}}n\right).\label{eq:criticalprob}
\end{align}
The same conclusions hold with $\bar\la$ in place of $\la$, including
when the critical-window condition is imposed on $\bar\la$.
\end{theorem}
Thus the Poisson estimate has a strictly positive leading bias in the
critical window: the true avoidance probability is larger.

\begin{theorem}[Optimal error and its lattice oscillation]\label{thm:lattice}
For either parameter $\eta$ in Theorem~\ref{thm:uniform}, let
\[
 \mathcal E_n(\eta)=\sup_{2\le r\le n}
       |\PP(U_n<r)-e^{-\eta_{n,r}}|,
 \quad
 t_n=\frac{2\log n-\log\log n+\log(q\alpha/4)}{\alpha},
\]
and define the continuous one-periodic function
\begin{equation}
 \Phi_p(\theta)=\max_{j\in\mathbb Z}
 \left\{e^{2\alpha(\theta-j)}
              \exp\bigl(-e^{\alpha(\theta-j)}\bigr)\right\}.
 \label{eq:Phi}
\end{equation}
Then
\begin{equation}
 \mathcal E_n(\eta)=\frac{4c_p}{\alpha^2}\frac{(\log n)^2}{n}
              \Phi_p(\{t_n\})
       +O_p\!\left(\frac{(\log n)^{3/2}}n\right).
 \label{eq:latticeerror}
\end{equation}
In particular $\mathcal E_n(\eta)=\Theta_p((\log n)^2/n)$, and
\begin{align}
 \limsup_{n\to\infty}\frac{n\mathcal E_n(\eta)}{(\log n)^2}
   &=\frac{16c_p}{\alpha^2e^2},\label{eq:limsup}\\
 \liminf_{n\to\infty}\frac{n\mathcal E_n(\eta)}{(\log n)^2}
   &=\frac{16c_p}{q^2}\exp\left(-\frac{2\alpha}{q}\right).
   \label{eq:liminf}
\end{align}
These two constants are distinct for every fixed $p\in(0,1)$.
\end{theorem}

\subsection{What is inherited, and what is being added}
Benjamini, Yadin, and Zeitouni studied longest progressions in random
subsets using Poisson approximation~\cite{BYZ}.
Zhao and Zhang's Theorem~1.1(1), in the preprint inspected here, gives
\begin{equation}
 \sup_{2\le r\le n}|\PP(U_n<r)-e^{-\bar\la_{n,r}}|
 =O_p\!\left(\frac{\log^4 n\,\log\log n}{n}\right).
 \label{eq:ZZ}
\end{equation}
Their progression-head construction and overlap estimates are directly
relevant~\cite{ZZ}. We explicitly reuse and reprove those ingredients;
neither the construction nor the three-term pair estimate below is
claimed as new. Theorems~\ref{thm:uniform}--\ref{thm:lattice} refine the
displayed estimate~\eqref{eq:ZZ}, rather than merely substitute constants
in its proof.

Higher-order avoidance estimates through joint cumulants are an established
subject. In particular, Mousset, Noever, Panagiotou, and Samotij give
systematic expansions for binomial-subset nonexistence probabilities,
including fixed-length arithmetic progressions~\cite{MNPS}.
Our event family is nonmonotone because a progression head requires a
zero immediately before a string of ones. Section~\ref{sec:avoidance}
therefore proves the particular local expansion needed here directly.
We make no priority claim for the general principle that the second
factorial cumulant supplies the first Poisson correction.

The ProveIt Ramsey directory contains research branches devoted to Gowers'
proof, quasipolynomial progressions, square differences, and van der Waerden
problems~\cite{ProveIt}. The OpenAI collection includes an arithmetic-
progression manuscript family and explicitly distinguishes verification
statuses~\cite{OpenAIMath,OpenAIAP}. These motivate the choice of topic,
but are not premises of any theorem proved here. This paper belongs in
probabilistic Ramsey/additive combinatorics, not in the dependency chain
of a deterministic density-increment proof.

\section{Progression heads and the exact mean}\label{sec:heads}
Let
\[
 \AP_{n,r}=\{(a,d):a,d\in\mathbb Z_{>0},\ a+(r-1)d\le n\}.
\]
For $i=(a,d)$, define the positive core and support
\[
 C_i=\{a,a+d,\ldots,a+(r-1)d\},\qquad
 S_i=C_i\cup\bigl(\{a-d\}\cap[n]\bigr).
\]
The head event is
\begin{equation}
 E_i=\begin{cases}
       \{\xi_x=1\text{ for all }x\in C_i\},&a\le d,\\
       \{\xi_{a-d}=0\}\cap\{\xi_x=1\text{ for all }x\in C_i\},&a>d,
     \end{cases}
 \qquad Y=\sum_{i\in\AP_{n,r}}\ind_{E_i}.
 \label{eq:head}
\end{equation}
A boundary head has $a\le d$ and probability $p^r$; an interior head has
probability $qp^r$.

\begin{lemma}[Exact reduction]\label{lem:mean}
For every $n\ge2$ and $2\le r\le n$,
\[
 \{U_n<r\}=\{Y=0\},\qquad \EE Y=\la_{n,r}.
\]
Moreover $Y_{n,r+1}\le Y_{n,r}$ pointwise, and therefore $\la_{n,r}$
is nonincreasing in $r$. Uniformly in $2\le r\le n$,
\begin{equation}
 |\la_{n,r}-\bar\la_{n,r}|\le 2np^r.\label{eq:meanshift}
\end{equation}
If $r=O(\log n)$ and $r\to\infty$, then
\begin{equation}
 \la_{n,r}=\frac{qn^2p^r}{2r}
          \left(1+O_p\left(\frac1r+\frac rn\right)\right).
 \label{eq:meanasy}
\end{equation}
\end{lemma}
\begin{proof}
Any selected $r$-term progression can be extended backwards along its
common difference until the first selected element of that run is reached.
The resulting run begins at a head and still has length at least $r$.
Conversely every head supplies an $r$-term progression.
The same starting point and difference of an $(r+1)$-head also define an
$r$-head, proving pointwise monotonicity.

There are $M_r(n)$ possible cores. The interior ones correspond bijectively
to $(r+1)$-term progressions by adjoining the predecessor. Replacing their
probability $p^r$ by $qp^r$ proves~\eqref{eq:lambda}.
For $m=r-1$, put $u=n/m$ and $t=D_r$ in~\eqref{eq:M}. We have
$0<u-t\le1$, and a direct expansion gives
\[
 M_r(n)=\frac{n^2}{2m}-\frac n2
       +\frac m2\bigl((u-t)-(u-t)^2\bigr).
\]
Thus $|M_r(n)-n^2/(2(r-1))|\le n$; the same bound holds for
$M_{r+1}$. This proves~\eqref{eq:meanshift}. Combining it with the exact
formula~\eqref{eq:barlambda} proves~\eqref{eq:meanasy}.
\end{proof}

\section{A local expansion for avoidance of cylinder events}\label{sec:avoidance}
This section is independent of arithmetic progressions. Its purpose is to
justify the passage from a covariance calculation to a zero-probability
asymptotic, with a quantified remainder. It does not assume positivity
of correlations or monotonicity of the events.

\subsection{Parameters and elementary conditional bounds}
Let $(E_i)_{i\in I}$ be a finite family of events, each determined by a
specified finite set of independent coordinates $S_i$. Join $i\ne j$
when $S_i\cap S_j\ne\varnothing$, write $i\sim j$, and let
$B_i=\{i\}\cup\{j:j\sim i\}$. Set
\begin{align}
 \pi_i&=\PP(E_i),\quad \pi_{ij}=\PP(E_i\cap E_j),\quad
 \pi_{ijk}=\PP(E_i\cap E_j\cap E_k),\nonumber\\
 \la&=\sum_i\pi_i,\quad
 \delta_i=\sum_{j\in B_i}\pi_j,\quad \delta=\max_i\delta_i,\nonumber\\
 b_1&=\sum_i\pi_i\delta_i,\qquad
 b_2=\sum_i\sum_{j\sim i}\pi_{ij},\label{eq:localpars}\\
 T&=\sum_{\substack{\{i,j,k\}\subset I\\
                 \text{induced graph connected}}}\pi_{ijk}.
 \label{eq:Tdef}
\end{align}
The pairs in $b_2$ are ordered; the triples in $T$ are unordered and
distinct. Write $F_A=\bigcap_{j\in A}E_j^c$.

\begin{lemma}[Conditional control]\label{lem:conditional}
Suppose $\delta\le1/8$. Every $F_A$ has positive probability, and for
$i\notin A$,
\begin{equation}
 \PP(E_i\mid F_A)\le2\pi_i.\label{eq:LLL}
\end{equation}
More generally, let $G$ be an event determined by the union of at most $m$
of the supports $S_i$. Let $N_G$ be the union of the closed neighborhoods of those at most $m$
supports and put $w=\sum_{j\in A\cap N_G}\pi_j$. Then
\begin{align}
 \PP(G\mid F_A)&\le\PP(G)e^{4w},\label{eq:auxupper}\\
 \bigl(\PP(G\mid F_A)-\PP(G)\bigr)_+
       &\le4e^{4m\delta}\PP(G)w,\label{eq:auxexcess}\\
 \bigl(\PP(G)-\PP(G\mid F_A)\bigr)_+
       &\le e^{4(m+1)\delta}
           \sum_{j\in A\cap N_G}\PP(G\cap E_j).
           \label{eq:auxdeficit}
\end{align}
\end{lemma}
\begin{proof}
We give the usual local-lemma induction to fix the conditional statement.
For $A=\varnothing$,~\eqref{eq:LLL} holds. Assume the conditional bound
for conditioning sets of size less than $|A|$. First obtain positivity of
$F_A$ by removing one index, using that bound, and observing that
$1-2\pi_i>0$. Then split $A$ into the neighbors
$V=A\cap B_i$ and the nonneighbors $W=A\setminus B_i$.
When $V=\varnothing$, independence gives the claim. Otherwise
\[
 \PP(E_i\mid F_A)
 \le \frac{\PP(E_i\mid F_W)}{\PP(F_V\mid F_W)}
 \le\frac{\pi_i}{\prod_{j\in V}(1-2\pi_j)}
 \le\frac{\pi_i}{1-2\sum_{j\in V}\pi_j}
 \le2\pi_i.
\]
The product bound uses the induction hypothesis, adding the elements of
$V$ one at a time; every intermediate conditioning set has smaller size.
Positivity at the next size follows by adding one event complement and
using $2\pi_i\le1/4<1$.

For $G$, split $A$ into $V=A\cap N_G$ and $W=A\setminus N_G$.
The event $G$ is independent of all coordinates used by $F_W$. The same
product argument and $-\log(1-2x)\le4x$ for $0\le x\le1/8$ give
\[
 \PP(G\mid F_A)
 \le\frac{\PP(G)}{\prod_{j\in V}(1-2\pi_j)}
 \le\PP(G)e^{4w}.
\]
Since $w\le m\delta$, $e^{4w}-1\le4we^{4m\delta}$ proves the upper
excess estimate. For the lower estimate, the denominator
$\PP(F_V\mid F_W)$ is at most one, so the union bound yields
\[
 \PP(G\mid F_A)\ge
 \PP(G)-\sum_{j\in V}\PP(G\cap E_j\mid F_W).
\]
Each intersection uses at most $m+1$ original supports. Applying the
just-proved upper estimate proves~\eqref{eq:auxdeficit}.
\end{proof}

\subsection{An exponentially damped first-order estimate}
\begin{lemma}[Coarse logarithmic sandwich]\label{lem:sandwich}
Under $\delta\le1/8$,
\begin{equation}
 -\la-4b_1\le\log\PP(F_I)\le-\la+2b_2.
 \label{eq:sandwich}
\end{equation}
If also $b_2\le\la/4$, then
\begin{equation}
 |\PP(F_I)-e^{-\la}|\le(4b_1+2b_2)e^{-\la/2}.
 \label{eq:damped}
\end{equation}
\end{lemma}
\begin{proof}
Order $I$, and put $z_i=\PP(E_i\mid F_{\{j:j<i\}})$ and
$u_i=\sum_{j<i,\,j\sim i}\pi_j$.
Lemma~\ref{lem:conditional} gives
\[
 z_i\le\pi_i e^{4u_i},\qquad
 z_i\ge\pi_i-e^{8\delta}\sum_{j<i,\,j\sim i}\pi_{ij}.
\]
The second inequality summed over $i$, followed by
$\log(1-z)\le-z$, gives the upper bound because
$e^{8\delta}/2\le e/2<2$.
For the lower bound,
\[
 \sum_i(z_i-\pi_i)
 \le4e^{4\delta}\sum_i\pi_i u_i
 \le4\left(b_1-\sum_i\pi_i^2\right).
\]
Here $\sum_i\pi_i u_i=(b_1-\sum_i\pi_i^2)/2$ and
$2e^{1/2}<4$. As $z_i\le2\pi_i\le1/4$, we may use
$\log(1-z_i)\ge-z_i-z_i^2$ and
$\sum_i z_i^2\le4\sum_i\pi_i^2$. This proves~\eqref{eq:sandwich}.
Finally write $\PP(F_I)=e^{-\la+v}$ with
$-4b_1\le v\le2b_2$. The elementary inequality
$|e^v-1|\le|v|e^{\max(v,0)}$ proves~\eqref{eq:damped}.
\end{proof}

\subsection{The second factorial cumulant, with a triple remainder}
\begin{proposition}[Second-order local avoidance]\label{prop:local}
There is an absolute constant $C$ such that, whenever $\delta\le1/8$,
\begin{equation}
 \left|\log\PP(F_I)+\la-
       \frac{\Var(\sum_i\ind_{E_i})-\la}{2}\right|
 \le C\bigl(\delta(b_1+b_2)+T\bigr).
 \label{eq:localsecond}
\end{equation}
\end{proposition}
\begin{proof}
We include the bookkeeping because an unexplained cumulant truncation
would not justify the main result. Use the same order as above, and set
\[
 V_i=\{j<i:j\sim i\},\qquad W_i=\{j<i:j\not\sim i\}.
\]
In this proof all conditional quantities carrying subscript $i$ are
conditioned on $F_{W_i}$. Define
\[
 a_i=\sum_{j\in V_i}\PP(E_j\mid F_{W_i}),\qquad
 s_i=\sum_{j\in V_i}\PP(E_iE_j\mid F_{W_i}).
\]
Bonferroni's inequalities and independence of $E_i$ from $F_{W_i}$ give
\begin{align*}
 \PP(E_i\cap F_{V_i}\mid F_{W_i})&=\pi_i-s_i+\rho_i,\\
 \PP(F_{V_i}\mid F_{W_i})&=1-a_i+\eta_i,
\end{align*}
where
\begin{align*}
 0\le\rho_i&\le e^{12\delta}
       \sum_{\{j,k\}\subset V_i}\pi_{ijk},\\
 0\le\eta_i&\le e^{8\delta}
       \sum_{\{j,k\}\subset V_i}\pi_{jk},\qquad \eta_i\le a_i.
\end{align*}
These bounds follow from~\eqref{eq:auxupper}. Also
$a_i\le2\delta_i\le1/4$, so the denominator is at least $3/4$.
Expanding this quotient, with no infinite series, gives
\begin{equation}
 z_i=\pi_i-s_i+\pi_i a_i+R_i,\qquad
 |R_i|\le C\bigl(\rho_i+a_i s_i+\pi_i a_i^2+\pi_i\eta_i\bigr).
 \label{eq:quotient}
\end{equation}
For completeness, subtract $\pi_i-s_i+\pi_i a_i$ from the quotient;
the numerator is
$\rho_i-s_i a_i-(\pi_i-s_i)\eta_i+\pi_i a_i^2-\pi_i a_i\eta_i$.
Use $\eta_i\le a_i\le1/4$ and divide by at least $3/4$.

We now sum each error. The triples in $\rho_i$ are connected stars, each
counted at most three times, hence $\sum_i\rho_i\le C T$.
Further,
\[
 \sum_i a_i s_i\le C\delta b_2,\qquad
 \sum_i\pi_i a_i^2\le4\delta b_1.
\]
For $\sum_i\pi_i\eta_i$, split the pair $j,k\in V_i$ according to
whether $j\sim k$. Nonadjacent pairs are independent, giving at most
$C\sum_i\pi_i\delta_i^2\le C\delta b_1$. For adjacent pairs, first
sum over $i\in B_j\cap B_k$: the sum of their $\pi_i$ is at most
$\delta$. This gives at most $C\delta b_2$. Consequently
\begin{equation}
 \sum_i|R_i|\le C\bigl(T+\delta(b_1+b_2)\bigr).
 \label{eq:sumR}
\end{equation}

Next replace the conditional pair sum $s_i$ by
$v_i=\sum_{j\in V_i}\pi_{ij}$.
The excess estimate~\eqref{eq:auxexcess}, applied to $E_iE_j$, costs at
most $C\delta\pi_{ij}$. The deficit estimate~\eqref{eq:auxdeficit}
costs at most $C\sum_{k\in W_i\cap B_j}\pi_{ijk}$:
$k$ cannot meet $S_i$ because it belongs to $W_i$. These are connected
three-vertex paths, and any unordered triple is counted at most six times.
It follows that
\begin{equation}
 \sum_i|s_i-v_i|\le C(\delta b_2+T).\label{eq:sreplace}
\end{equation}
Similarly replace $a_i$ by $u_i=\sum_{j\in V_i}\pi_j$, remembering the
factor $\pi_i$. The excess is bounded by
$C\delta\sum_i\pi_i\sum_{j\in V_i}\pi_j\le C\delta b_1$.
The deficit is bounded by
\[
 C\sum_i\pi_i\sum_{j\in V_i}\sum_{k\in W_i\cap B_j}\pi_{jk}.
\]
After summing over $i\in B_j$, this is at most $C\delta b_2$.
Together with~\eqref{eq:quotient}--\eqref{eq:sreplace}, this proves
\begin{equation}
 \sum_i\left|z_i-\pi_i-
       \sum_{j\in V_i}(\pi_i\pi_j-\pi_{ij})\right|
 \le C\bigl(\delta(b_1+b_2)+T\bigr).
 \label{eq:zexpansion}
\end{equation}

It remains to pass from conditional probabilities to their logarithms.
The first-order bounds in Lemma~\ref{lem:conditional} imply
\[
 \sum_i|z_i-\pi_i|\le C(b_1+b_2).
\]
Since $z_i+\pi_i\le3\pi_i\le3\delta$,
$\sum_i|z_i^2-\pi_i^2|\le C\delta(b_1+b_2)$.
Also $\sum_i z_i^3\le8\delta\sum_i\pi_i^2\le8\delta b_1$.
Taylor's formula on $[0,1/4]$ now gives
\begin{align*}
 \log\PP(F_I)
 &=\sum_i\log(1-z_i)\\
 &=-\la+\sum_{i}\sum_{j\in V_i}(\pi_{ij}-\pi_i\pi_j)
       -\frac12\sum_i\pi_i^2
       +O\bigl(\delta(b_1+b_2)+T\bigr).
\end{align*}
The middle two terms are $(\Var(\sum_i\ind_{E_i})-\la)/2$:
all covariances of nonadjacent events vanish. This proves the proposition.
\end{proof}

\section{Geometry of pairs and connected triples}\label{sec:geometry}
Return to progression heads, with dependency graph defined by their
supports. Throughout this section $r\ge3$; absolute implied constants
are sufficient. Put $\pi_i=\PP(E_i)\le p^r$.

\subsection{Pair counts and compatible overlaps}
\begin{lemma}[Pair geometry; compare Zhao--Zhang, Lemma 3.2]\label{lem:pairs}
The number of ordered support-intersecting pairs is $O(n^3)$.
The number with at least two shared support sites is $O(n^2r^3)$.
Among intersecting pairs with common differences in ratio two the number
is $O(n^2)$. Every head has at most $O(nr)$ neighbors, and
\begin{align}
 \delta&\ll nrp^r,\label{eq:deltaAP}\\
 b_1+b_2&\ll n^3p^{2r-1}
        +n^2r^3p^{5r/3-1}+n^2p^{3r/2-1}.
        \label{eq:pairbound}
\end{align}
For a compatible pair of distinct heads with different common differences
$d,e$, put $m=\max(d/g,e/g)$, where $g=\gcd(d,e)$. Then
\begin{equation}
 |S_i\cap S_j|\le r/m+1.\label{eq:overlap}
\end{equation}
Distinct, intersecting heads with the same common difference are
incompatible.
\end{lemma}
\begin{proof}
A fixed site belongs to at most
$(r+1)\lfloor(n-1)/(r-1)\rfloor=O(n)$ supports: choose its index in
$\{-1,0,\ldots,r-1\}$ and the common difference. As a support has at
most $r+1$ sites, this gives $O(nr)$ neighbors. Since
$|\AP_{n,r}|=O(n^2/r)$, there are $O(n^3)$ ordered intersecting pairs.

If two supports share two sites, their common differences have ratio
$u/v$ for integers $1\le u,v\le r$. Once the first head is fixed,
there are at most $r^2$ choices of this ratio and at most $(r+1)^2$
choices of indices specifying a common point. Hence there are
$O(r^4)$ such second heads and $O(n^2r^3)$ such pairs.
For $e=2d$, a common point forces $b=a+kd$ with $k$ in an interval
of length $O(r)$; there are only $O(r)$ second starts for each first
head. Interchanging $d,e$ proves the $O(n^2)$ estimate.

Two common support points differ by a multiple of $\operatorname{lcm}(d,e)$.
Their indices in the two supports therefore differ by multiples of
$e/g$ and $d/g$, respectively. The available index spans are at most $r$,
which proves~\eqref{eq:overlap}.
When $d=e$ and $b>a$, an intersection gives $b=a+kd$ for
$1\le k\le r$. The later head requires $\xi_{b-d}=0$, while that site
belongs to the earlier positive core. The intersection event is empty.

A one-site-overlap pair has at least $2r-1$ distinct required positive
sites, and its joint probability is at most $p^{2r-1}$. For different
differences in ratio two it is at most $p^{3r/2-1}$. For all other
compatible multiply-overlapping pairs, $m\ge3$ and the joint probability
is at most $p^{5r/3-1}$. Multiplying these bounds by the counts proves the
$b_2$ estimate. The $b_1$ estimate $O(n^3p^{2r})$ is absorbed by the
first term, and~\eqref{eq:deltaAP} follows from the degree bound.
\end{proof}

\subsection{A triple estimate adequate for a second-order law}
\begin{lemma}[Connected triples]\label{lem:triples}
For the $T$ in~\eqref{eq:Tdef},
\begin{equation}
 T\ll n^4r^4p^{3r-3}
       +n^3r^5p^{5r/2-3}
       +n^2r^8p^{5r/3-3}.\label{eq:triplebound}
\end{equation}
\end{lemma}
\begin{proof}
A triple contributes only if the three head requirements are compatible.
Classify its pair edges as multiple when at least two support sites are
shared. There are three cases.

\emph{No multiple edges.} Choose a root head and then attach two heads
along a spanning tree. There are $O(n^2)$ choices for the root and
$O(nr)$ choices for each attachment, giving at most $O(n^4r^4)$ triples
(the powers used here are deliberately loose). Each pair of positive
cores intersects in at most one site, so the union of the three cores
has at least $3r-3$ sites.

\emph{Exactly one multiple edge.} Choose that pair in $O(n^2r^3)$ ways.
The third head meets at least one member, giving $O(nr)$ choices, hence
at most $O(n^3r^5)$ triples. The multiple pair, if compatible, has
intersection at most $r/2+1$; each other intersection has size at most
one. The positive union has size at least $5r/2-3$.

\emph{At least two multiple edges.} Two multiple edges contain a spanning
tree. Starting from a root, each attachment has $O(r^4)$ choices by the
two-site argument. There are at most $O(n^2r^8)$ triples.
If all three common differences are distinct, every pair overlap is at
most $r/2+1$, and at least one is at most $r/3+1$. Indeed three distinct
positive integers cannot have every pair in ratio two. Inclusion--exclusion
therefore bounds the positive union below by $5r/3-3$.
If two common differences coincide, compatibility forces their supports
to be disjoint. The other two overlaps are at most $r/2+1$, giving the
stronger lower bound $2r-2$. If all coincide, connectedness and
compatibility cannot both hold.

In each case ignore required zeros, multiply the count by $p$ to the
appropriate lower bound on the positive union, and sum. This gives the
stated estimate even when a displayed fractional exponent is negative
for small $r$, since such an upper bound remains valid.
\end{proof}

\section{The entropy profile of progression incidence}\label{sec:entropy}
For $x\in[n]$, let $d_r(x)$ be the number of $r$-term positive cores
containing $x$. The main coefficient comes from these degrees, not from
the rare high-multiplicity intersections.

\begin{lemma}[Uniform degree profile]\label{lem:degree}
Let $h(u)=-u\log u-(1-u)\log(1-u)$, with $h(0)=h(1)=0$.
Uniformly in $x\in[n]$, for $3\le r\le n$,
\begin{equation}
 d_r(x)=n h\left(\frac{x-1}{n-1}\right)
              +O\left(\frac nr+r\right).\label{eq:degreeprofile}
\end{equation}
Consequently, uniformly for $3\le r\le n$,
\begin{equation}
 \sum_{x=1}^n d_r(x)^2
       =n^3H+O\left(\frac{n^3}{r}+n^2r\right).
 \label{eq:degreesquare}
\end{equation}
\end{lemma}
\begin{proof}
Put $m=r-1$ and $u=(x-1)/(n-1)$. If $x$ is at position $j$ in its
progression, $0\le j\le m$, the permissible differences are the positive
integers at most
\[
 \min\left(\frac{x-1}{j},\frac{n-x}{m-j}\right),
\]
where an endpoint's missing constraint is omitted. Thus
\[
 d_r(x)=\frac{n-1}{m}\sum_{j=0}^{m}f_u(j/m)+O(r),
 \quad
 f_u(t)=\min\left(\frac ut,\frac{1-u}{1-t}\right),
\]
with $f_u(0)=1-u$ and $f_u(1)=u$. The function increases up to its peak
at $t=u$ and then decreases, has height at most one, and has total
variation at most two, including $u=0,1$. Therefore its Riemann sum
satisfies
\[
 \frac1m\sum_{j=0}^m f_u(j/m)=\int_0^1 f_u(t)\dd t+O(1/m)
\]
uniformly in $u$. Splitting the integral at $u$ gives
\[
 \int_0^u\frac{1-u}{1-t}\dd t
 +\int_u^1\frac u t\dd t=h(u).
\]
This proves~\eqref{eq:degreeprofile}.

The function $h^2$ is bounded and has bounded variation, so its Riemann
sum on the $n$ grid points is $nH+O(1)$. Squaring the degree formula
and summing contributes an error at most
\[
 O\left(n^2\left(\frac nr+r\right)
          +n\left(\frac nr+r\right)^2+n^2\right).
\]
For $3\le r\le n$ this is
$O(n^3/r+n^2r)$, proving~\eqref{eq:degreesquare}. In particular,
$n^{-3}\sum_xd_r(x)^2\to H$ whenever $r\to\infty$ and $r=o(n)$.
\end{proof}

\begin{lemma}[Evaluation of the entropy integral]\label{lem:H}
The value of $H$ is $5/6-\pi^2/18$.
\end{lemma}
\begin{proof}
Expansion of the square and symmetry give
\[
 H=2\int_0^1u^2\log^2u\dd u
       +2\int_0^1u(1-u)\log u\log(1-u)\dd u.
\]
The first integral is $2/27$. Expanding $-\log(1-u)$ in its nonnegative
power series and integrating by monotone convergence evaluates the second
integral as
\[
 J=\sum_{k\ge1}\frac1k
       \left(\frac1{(k+2)^2}-\frac1{(k+3)^2}\right).
\]
For an integer $a\ge1$, partial fractions give
\[
 \sum_{k\ge1}\frac1{k(k+a)^2}
   =\frac{H_a}{a^2}-\frac{\zeta(2)-H_a^{(2)}}a,
\]
where $H_a=\sum_{j=1}^a1/j$ and $H_a^{(2)}=\sum_{j=1}^a1/j^2$.
Taking the case $a=2$ minus the case $a=3$ yields
$J=37/108-\pi^2/36$. Hence
$H=4/27+2J=5/6-\pi^2/18$.
\end{proof}

\section{Computing the leading covariance term}\label{sec:variance}
For a head $i$ and a coordinate $x$, define its signed influence by
\[
 D_{i,x}=\EE(\ind_{E_i}\mid\xi_x=1)
           -\EE(\ind_{E_i}\mid\xi_x=0).
\]
It is zero outside $S_i$. If $x$ belongs to the positive core it equals
$p^{r-1}$ for a boundary head and $qp^{r-1}$ for an interior head.
At the predecessor of an interior head it equals $-p^r$.

\begin{lemma}[One-site covariance identity]\label{lem:influence}
If $S_i\cap S_j=\{x\}$, then
\begin{equation}
 \Cov(\ind_{E_i},\ind_{E_j})=pqD_{i,x}D_{j,x}.\label{eq:onecov}
\end{equation}
Let $B_r(x)$ count boundary heads whose positive core contains $x$,
and put $P_r(x)=\lfloor(n-x)/r\rfloor$. Then, exactly,
\begin{equation}
 S(x):=\sum_iD_{i,x}
   =qp^{r-1}d_r(x)+p^rB_r(x)-p^rP_r(x),\qquad
 B_r(x),P_r(x)\le \frac{n}{r-1}.\label{eq:Stotal}
\end{equation}
\end{lemma}
\begin{proof}
Given $\xi_x$, the remaining coordinates of the two events are disjoint
and hence independent. The covariance of two affine functions of a
Bernoulli-$p$ variable is $pq$ times the product of their slopes, proving
\eqref{eq:onecov}.
The formula for $S(x)$ follows from the three influence values above.
A head with predecessor $x$ has start $x+d$ and endpoint $x+rd\le n$,
so there are exactly $P_r(x)$ such heads. For each fixed common difference
$d$, at most one boundary start $a\in\{1,\ldots,d\}$ can have $x$ in its
core: it is the unique representative of $x$ modulo $d$ in that interval.
This bounds $B_r(x)$ by $D_r$.
\end{proof}

\begin{proposition}[Second factorial cumulant]\label{prop:variance}
Fix $p\in(0,1)$, and suppose $r\to\infty$ with $r=O(\log n)$.
Then
\begin{align}
 \frac{\Var Y-\la}{2}
  &=\frac{q^3H}{2p}\,n^3p^{2r}+O_p(V_{n,r}),\label{eq:variancelead}\\
 V_{n,r}
  &=\frac{n^3p^{2r}}r+n^2rp^{2r}
    +n^2r^4p^{2r-2}
    +n^2p^{3r/2-1}+n^2r^3p^{5r/3-1}.
 \label{eq:V}
\end{align}
\end{proposition}
\begin{proof}
By~\eqref{eq:Stotal}, $S(x)=qp^{r-1}d_r(x)+O(p^rn/r)$.
Using Lemma~\ref{lem:degree} and $d_r(x)=O(n)$ yields
\begin{equation}
 \frac{pq}{2}\sum_x S(x)^2
 =\frac{q^3H}{2p}n^3p^{2r}
       +O_p\left(\frac{n^3p^{2r}}r+n^2rp^{2r}\right).
 \label{eq:projectionnorm}
\end{equation}
Expanding the square on the left gives a diagonal sum over individual
heads plus sums over unordered pairs. By~\eqref{eq:onecov}, the
one-site-overlap pairs have exactly their true covariance contribution.
The diagonal projection has size $O_p(n^2p^{2r-2})$: there are
$O(n^2/r)$ heads and $O(r)$ nonzero influences per head.
The projection contributions of multiply-overlapping pairs are at most
\[
 O_p\bigl((n^2r^3)\,r\,p^{2r-2}\bigr)
       =O_p(n^2r^4p^{2r-2}).
\]
Their actual covariance sum in absolute value is bounded by their joint
probability sum plus the product-probability sum. Lemma~\ref{lem:pairs}
therefore bounds it by
\[
 O\left(n^2p^{3r/2-1}+n^2r^3p^{5r/3-1}
                     +n^2r^3p^{2r}\right).
\]
Finally, the second factorial cumulant subtracts
$\frac12\sum_i\pi_i^2=O(n^2p^{2r}/r)$ from the unordered covariance sum.
All diagonal and product terms are absorbed by the third term of
\eqref{eq:V}. Combining these observations proves the proposition.
\end{proof}

\begin{remark}[Why the sign is positive]
The dominant term is a sum of squares of aggregate signed influences.
Although predecessors introduce negative influences and same-difference
heads can exclude one another, their aggregate contribution is smaller
by a factor of at least $1/r$ in~\eqref{eq:Stotal}. The leading term is
therefore the positive entropy integral. Positivity is a conclusion of
this calculation, not an application of positive association to the
nonmonotone head events.
\end{remark}

\section{Assembly of the critical and uniform expansions}\label{sec:assembly}
\subsection{A logarithmic band with a weighted remainder}
\begin{lemma}[Uniform expansion on a logarithmic band]\label{lem:band}
Fix positive constants $c_1,c_2,C_0$. Suppose
\[
 c_1\log n\le r\le c_2\log n,\qquad
 0<\la=\la_{n,r}\le C_0\log n.
\]
Then, for all sufficiently large $n$,
\begin{equation}
 \log\PP(U_n<r)
 =-\la+c_p\la^2\frac{r^2}{n}+R_{n,r},\qquad
 |R_{n,r}|\ll_{p,c_1,c_2,C_0}
       (\la^{3/2}+\la^3)\frac{r^{3/2}}n.
 \label{eq:bandlog}
\end{equation}
\end{lemma}
\begin{proof}
On this band~\eqref{eq:meanasy} gives
$p^r\asymp_p\la r/n^2$. In particular
$\delta\ll_p\la r^2/n=O((\log n)^3/n)$, so the local avoidance
lemmas apply. The pair and triple estimates become
\begin{align}
 b_1+b_2&\ll_p
       \la^2\frac{r^2}{n}
       +\la^{5/3}\frac{r^{14/3}}{n^{4/3}}
       +\la^{3/2}\frac{r^{3/2}}n,\label{eq:scaledpairs}\\
 T&\ll_p
       \la^3\frac{r^7}{n^2}
       +\la^{5/2}\frac{r^{15/2}}{n^2}
       +\la^{5/3}\frac{r^{29/3}}{n^{4/3}}.\label{eq:scaledtriples}
\end{align}
Also~\eqref{eq:meanasy} converts the main covariance term into
\begin{equation}
 \frac{q^3H}{2p}n^3p^{2r}
 =c_p\la^2\frac{r^2}n
       \left(1+O_p\left(\frac1r+\frac rn\right)\right).
 \label{eq:convertconstant}
\end{equation}
Substituting into $V_{n,r}$ in~\eqref{eq:V} produces terms bounded by
\[
 O_p\left(
  \la^2\frac rn+\la^2\frac{r^6}{n^2}
  +\la^{3/2}\frac{r^{3/2}}n
  +\la^{5/3}\frac{r^{14/3}}{n^{4/3}}\right).
\]
The error in~\eqref{eq:convertconstant} is absorbed in these terms.
Combining Proposition~\ref{prop:local} with
Proposition~\ref{prop:variance}, including the product
$\delta(b_1+b_2)$ from~\eqref{eq:scaledpairs}, bounds the total
remainder by
\begin{equation}
 C_p(\la^{3/2}+\la^3)
 \left(\frac{r^{3/2}}n+\frac{r^{10}}{n^{4/3}}
                               +\frac{r^8}{n^2}\right).
 \label{eq:unabsorbed}
\end{equation}
To see that the powers of $\la$ suffice, every exponent arising here is
between $3/2$ and $3$; $\la^s\le\la^{3/2}+\la^3$ in that range.
Since $r\asymp\log n$, the last two terms in parentheses are
$o(r^{3/2}/n)$, uniformly on the band. This proves~\eqref{eq:bandlog}.
\end{proof}

\begin{proof}[Proof of Theorem~\ref{thm:critical}]
First, $\la\asymp1$ forces $r\asymp\log n$.
Indeed the elementary upper bound $\la\le n^2p^r$ gives $r=O(\log n)$,
whereas $\la\ge qp^rM_r(n)$ and~\eqref{eq:M} give a matching lower
bound once $r=O(\log n)$. Fixed $r$ cannot keep $\la$ bounded as
$n\to\infty$. Taking logarithms in~\eqref{eq:meanasy} then yields
\eqref{eq:criticalr}.
Lemma~\ref{lem:band} proves~\eqref{eq:criticallog}.
Exponentiation proves~\eqref{eq:criticalprob}, because
$r^4/n^2=o(r^{3/2}/n)$. Finally~\eqref{eq:meanshift} is $O_{p,a,b}(r/n)$
in this window, smaller than the stated remainder. The same relative
mean estimate shows that a critical condition on $\bar\la$ is equivalent
to a slightly enlarged critical condition on $\la$.
\end{proof}

\subsection{Why the tails do not lose logarithms}
\begin{proof}[Proof of Theorem~\ref{thm:uniform} for the exact mean]
Put $L=\log n$. Let $r_-$ be the smallest $r\ge3$ with
$\la_{n,r}\le(16/p)L$, and let
$r_+=\lceil3L/\alpha\rceil$. For large $n$, both are of order $L$,
$r_-<r_+$, and
\begin{equation}
 8L\le\la_{n,r_-}\le(16/p)L.\label{eq:cutoff}
\end{equation}
To justify the lower inequality, uniformly for $r\asymp L$ the explicit
mean formula gives
$\la_{n,r}/\la_{n,r-1}=p(1+O_p(1/r+r/n))\ge p/2$.
The defining previous mean is larger than $(16/p)L$.

For $r_-\le r\le r_+$, the hypotheses of Lemma~\ref{lem:band} hold.
Its correction $K=c_p\la^2r^2/n$ and remainder $R$ are both uniformly
$o(1)$, since $\la=O_p(L)$. Hence
\begin{align*}
 \left|\PP(U_n<r)-e^{-\la}(1+K)\right|
 &\ll_p e^{-\la}\left((\la^{3/2}+\la^3)\frac{r^{3/2}}n
                               +\la^4\frac{r^4}{n^2}\right)\\
 &\ll_p \frac{L^{3/2}}n.
\end{align*}
The last inequality uses boundedness of $x^s e^{-x}$ on $[0,\infty)$.
This exponential damping is precisely what is lost in a uniform proof
that takes the worst overlap bound before controlling the avoidance tail.

For $r<r_-$,~\eqref{eq:scaledpairs} implies
$b_2=o(\la)$ at $r_-$.
Lemma~\ref{lem:sandwich} and~\eqref{eq:cutoff} give
$\PP(U_n<r_-)\le e^{-\la_{n,r_-}/2}\le n^{-4}$.
Both $\PP(U_n<r)$ and the uncorrected exponential are then negligible
by monotonicity. The correction is negligible as well: here
$r\le r_-=O_p(L)$ and $\la\ge8L$, so
$e^{-\la}\la^2r^2/n=O_p(L^4n^{-9})$.

For $r>r_+$, Markov's inequality and Lemma~\ref{lem:mean} give
$\PP(U_n\ge r)\le\la_{n,r}\le\la_{n,r_+}=O_p(1/(nL))$.
Also $1-e^{-\la}\le\la$. Finally the global bound
$\la\ll n^2p^r/r$ gives
\[
 e^{-\la}\la^2r^2/n\ll n^3p^{2r}\le n^{-3}.
\]
These estimates settle the upper tail and complete the exact-mean proof.
\end{proof}

\begin{proof}[Transfer to the smooth mean]
On the middle band,~\eqref{eq:meanshift} gives
$|\bar\la-\la|\ll_p\la r/n$ and
$\bar\la/\la=1+O_p(r/n)$. Apply the mean value theorem to
$g_r(x)=e^{-x}(1+c_px^2r^2/n)$. Its derivative is
\[
 g_r'(x)=e^{-x}\left[-1+c_p\frac{r^2}n(2x-x^2)\right].
\]
On the segment between the two means, the exponential is at most
$e^{-\la/2}$ for sufficiently large $n$. The resulting difference is
$O_p(L/n)$ uniformly, by exponential damping.
In the lower tail the relative mean estimate is still valid for
$2\le r<r_-$, using~\eqref{eq:M} and $r=O_p(L)$; thus
$\bar\la\ge\la/2\ge4L$. Both smooth terms are negligible.
In the upper tail $\bar\la\ll_p n^2p^r/r$, so the same estimates as above
apply. This proves~\eqref{eq:uniformcorrected} for $\bar\la$.
\end{proof}

\section{The optimal error has an explicit periodic amplitude}\label{sec:lattice}
\begin{proof}[Proof of Theorem~\ref{thm:lattice}]
Theorem~\ref{thm:uniform}, and the elementary inequality
$||x|-|y||\le|x-y|$, imply
\begin{equation}
 \mathcal E_n(\la)=\frac{c_p}{n}
       \max_{2\le r\le n}\left(e^{-\la_{n,r}}\la_{n,r}^2r^2\right)
             +O_p(L^{3/2}/n).\label{eq:errorasmax}
\end{equation}
Let $f(x)=x^2e^{-x}$. For $r=\lfloor t_n\rfloor+j$ with $j$ in any
fixed finite set,~\eqref{eq:meanasy} gives, uniformly,
\begin{equation}
 \la_{n,r}=e^{\alpha(\{t_n\}-j)}
       \left(1+O_p\left(\frac{\log L}{L}\right)\right),\qquad
 r^2=\frac{4L^2}{\alpha^2}+O_p(L\log L).\label{eq:latticemean}
\end{equation}
The constant in $t_n$ is determined by
$qn^2e^{-\alpha t_n}/(2(2L/\alpha))=1$.

We justify restricting the maximum to a fixed finite set of such $j$.
A lattice point with its limiting mean in $[p,1]$ provides
$f(\la)\ge m_p>0$ for all sufficiently large $n$. The two tails
$r<r_-$ and $r>r_+$ in the preceding proof contribute $o(L^2)$ to the
maximum in~\eqref{eq:errorasmax}. On the middle band $r^2/L^2$ is bounded
above and below by positive constants. Since $f(x)\to0$ at zero and
infinity, one can choose $0<a_p<b_p<\infty$ so that a maximizer must
have $a_p\le\la_{n,r}\le b_p$. Formula~\eqref{eq:criticalr}, with its
more precise constant supplied by~\eqref{eq:meanasy}, then gives
$r-t_n=O_p(1)$. The analogous maximum defining $\Phi_p$ also lies in a
fixed finite set of $j$, uniformly for $\theta\in[0,1]$.

The derivative of $f$ is bounded on the relevant compact interval.
Thus~\eqref{eq:latticemean} converts~\eqref{eq:errorasmax} to
\[
 \mathcal E_n(\la)=\frac{4c_pL^2}{\alpha^2n}\Phi_p(\{t_n\})
       +O_p\left(\frac{L\log L}{n}+\frac{L^{3/2}}n\right).
\]
Since $L\log L=O(L^{3/2})$, this is~\eqref{eq:latticeerror}.
The smooth parameter obeys the same lattice-mean expansion, so the same
argument applies directly to $\mathcal E_n(\bar\la)$.

It remains to evaluate the extrema of $\Phi_p$.
The unrestricted maximum of $f$ is $f(2)=4/e^2$, so
$\max_\theta\Phi_p(\theta)=4/e^2$.
For a multiplicative lattice of ratio $p$, only the two adjacent points
$x$ and $px$ straddling $2$ can maximize $f$. Thus $2\le x\le2/p$,
and its maximum is $\max(f(x),f(px))$.
The first term decreases and the second increases on that interval.
Their equality occurs at
\[
 x_*=-\frac{2\log p}{1-p}=\frac{2\alpha}{q}\in(2,2/p).
\]
Therefore
\begin{equation}
 \min_\theta\Phi_p(\theta)
       =\frac{4\alpha^2}{q^2}e^{-2\alpha/q}>0,
 \qquad \max_\theta\Phi_p(\theta)=\frac4{e^2}.
 \label{eq:Phiextrema}
\end{equation}
The maximum over a uniformly finite set also proves continuity of
$\Phi_p$, and relabeling $j$ proves periodicity.

Finally $t_n\to\infty$, is eventually increasing, and
$t_{n+1}-t_n\to0$. For every $\theta\in[0,1]$, take the first integer
$n$ crossing each successive level $k+\theta$. The overshoot tends to
zero, so the fractional parts $\{t_n\}$ have $\theta$ as an accumulation
point (with the endpoints understood periodically).
Equations~\eqref{eq:latticeerror} and~\eqref{eq:Phiextrema} now prove
\eqref{eq:limsup} and~\eqref{eq:liminf}. They are distinct because
$x_*>2$ and $f$ is strictly decreasing on $(2,\infty)$.
\end{proof}

\begin{remark}[Fair-coin specialization]
At $p=q=1/2$, the normalized error has
\[
 \liminf\frac{n\mathcal E_n}{(\log n)^2}=4c_{1/2}
       \approx2.280175821738063,
 \qquad
 \limsup\frac{n\mathcal E_n}{(\log n)^2}
       =\frac{16c_{1/2}}{(\log2)^2e^2}.
\]
The oscillation is not a numerical fitting phenomenon. It comes from the
integer spacing of progression lengths and the geometric spacing of their
Poisson means.
\end{remark}

\section{A separate sparse-density consequence}\label{sec:sparse}
The preceding main theorems fix $p$ and let $r$ grow. The head viewpoint
also gives an elementary consequence in the opposite regime, useful for
cross-checking what declumping does and does not remove.

\begin{proposition}[Fixed length, sparse density]\label{prop:sparse}
Fix $k\ge3$ and $C_0<\infty$. Uniformly for
$0\le c\le C_0$ and $p=c n^{-2/k}$ (for all sufficiently large $n$),
\begin{equation}
 \PP(U_n<k)
 =\exp\bigl(-M_k(n)p^k+M_{k+1}(n)p^{k+1}\bigr)
       +O_{k,C_0}\bigl(n^{-1+2/k}\bigr).
 \label{eq:sparse}
\end{equation}
For each fixed $c>0$ and $k\ge5$, this implies
\begin{equation}
 \log\PP(U_n<k)
 =-\frac{c^k}{2(k-1)}
       +\frac{c^{k+1}}{2k}n^{-2/k}
       +O_{k,c}\bigl(n^{-1+2/k}\bigr).
 \label{eq:sparselog}
\end{equation}
\end{proposition}
\begin{proof}
With $k$ fixed, there are $O_k(n^3)$ intersecting support pairs and
$O_k(n^2)$ multiply-overlapping pairs. A pair of cores with different
common differences shares at most $\lceil k/2\rceil$ positive sites:
the common index increments have maximum at least two. Same-difference
intersecting heads are incompatible. Thus
\[
 b_1\ll_k n^3p^{2k},\qquad
 b_2\ll_k n^3p^{2k-1}+n^2p^{2k-\lceil k/2\rceil}.
\]
At $p=c n^{-2/k}$ these are
$O_{k,C_0}(n^{-1+2/k})$; for the second term the exponents are $-1$
for even $k$ and $-1+1/k$ for odd $k$, both smaller.
Also $\delta\ll_k np^k=O_{k,C_0}(n^{-1})$.
Lemma~\ref{lem:sandwich} gives the claimed approximation, without
requiring a lower bound on $c$; the case $c=0$ is exact.
For fixed $c>0$ the avoidance probability stays bounded away from zero.
Taking logarithms and using $M_j(n)=n^2/(2(j-1))+O_j(n)$ proves
\eqref{eq:sparselog}. For $k\ge5$ its displayed correction exceeds
the remainder, since $2/k<1-2/k$.
\end{proof}

\begin{remark}[Interpretation, not a priority claim]
The exact mean incorporates the first sliding-run correction
$M_{k+1}p^{k+1}$. This is an illustration of declumping, not a claim that
fixed-pattern avoidance expansions originate here; the more general
fixed-length cumulant literature includes~\cite{MNPS}. The genuinely
separate issue in the main theorems is uniformity as $r\asymp\log n$,
together with evaluation and sharpness of the resulting leading error.
\end{remark}

\section{Verification, limitations, and integration}\label{sec:verification}
\subsection{What the executable artifact actually checks}
The supplied \texttt{code/verify.py} uses only Python's standard library.
It enumerates Bernoulli configurations and uses integers and rational
numbers for the finite event identities. The recorded full run passed:
\begin{center}
\begin{tabular}{lr}
\toprule
Check family & Number checked\\
\midrule
Configuration/threshold instances & 114,176\\
Exhaustive $(n,r)$ cases & 62\\
Pair parameter cases, including three rational densities & 36\\
One-site covariance identities & 21,183\\
Multiply-overlapping pair classifications & 9,546\\
Three-head geometric configurations & 810,589\\
Compatible connected triples among these & 194,077\\
Numerical local-avoidance sandwiches & 8\\
\bottomrule
\end{tabular}
\end{center}
The configuration checks verify the exact zero-event reduction, the mean
formula, and selected exact second moments. Pair checks also verify the
signed-influence sum and the projection decomposition. Triple checks
verify the three positive-union lower bounds used in
Lemma~\ref{lem:triples}. The eight logarithmic sandwich checks use
floating-point logarithms of exact rational probabilities and are labeled
numerical, rather than exact inequalities.

The degree-profile calculation checks the exact identity
$\sum_xd_r(x)=rM_r(n)$ and records
$n^{-3}\sum_xd_r(x)^2$ for $n$ up to $100{,}000$.
Its agreement with the entropy limit is a diagnostic, not evidence
sufficient to establish an asymptotic theorem.
The exact configuration enumeration is intentionally small. It neither
estimates a tiny large-$n$ correction by Monte Carlo nor silently treats
finite checks as a proof for all $n$.

\subsection{Claims that are not being made}
No theorem in this paper is currently certified by a proof assistant.
The manuscript's written proofs, especially
Proposition~\ref{prop:local} and its weighted use in
Lemma~\ref{lem:band}, require independent mathematical review.
The novelty comparison is specific to the inspected sources and is not
an exhaustive priority determination.
The remainder $O((\log n)^{3/2}/n)$ after correction is an upper bound;
we do not prove it is optimal.
No uniformity as $p$ approaches zero or one follows from the fixed-$p$
theorems. No claim is made for cyclic progressions modulo $n$, for a
fixed-cardinality random set, or for an improvement of deterministic
Ramsey numbers.

\subsection{Repository placement and proof dependencies}
A suggested destination, not an assertion that the folder already exists,
is
\begin{center}
\texttt{Combinatorics/Ramsey/Research/RandomProgressions/SharpPoisson/}.
\end{center}
The package contains this manuscript, its source, an external bibliography,
reproducible code, recorded data, a claims ledger, a literature/novelty
note, and a formalization roadmap. It does not modify either repository.
The proof dependencies are
\[
 \begin{gathered}
 \text{conditional control}\Longrightarrow\text{local avoidance expansion},\\[4pt]
 \left.\begin{matrix}
  \text{local avoidance expansion}\\
  \text{pair and triple geometry}\\
  \text{entropy covariance calculation}
 \end{matrix}\right\}
 \Longrightarrow\text{weighted band expansion},\\[4pt]
 \text{weighted band expansion}
 \Longrightarrow\text{uniform correction}
 \Longrightarrow\text{lattice error law}.
 \end{gathered}
\]

\section{Further research questions}\label{sec:questions}
The following questions are not used in the proofs above.

\paragraph{1. Evaluate the ratio-two correction.}
Pairs with common differences in ratio two account for the visible
$O(r^{3/2}/n)$ term in the critical window. Determine their exact
contribution, including dependence on the parity of $r$ and on the
lattice phase of $n$. Does subtracting it reduce the uniform remainder
to $O((\log n)/n)$? A solution would need actual covariance counts,
not just the upper bound in Lemma~\ref{lem:pairs}.

\paragraph{2. A third-order avoidance expansion with an effective constant.}
Extend Proposition~\ref{prop:local} by one order for nonmonotone cylinder
events. Separate the connected four-event terms from products of pair
terms, and supply explicit numerical constants valid under a stated
local-mass hypothesis. Such a result would support a systematic
expansion rather than an isolated correction.

\paragraph{3. Uniform regimes for varying $p=p_n$.}
Determine the largest ranges in which the entropy correction remains
dominant when $p_n\to0$ or $p_n\to1$. The present error terms contain
powers of $p^{-1}$ and $q^{-1}$ after mean normalization, so simply
replacing $p$ by $p_n$ in the main theorems is not justified. Locate the
transition between growing-pattern behavior and fixed-pattern sparse
avoidance.

\paragraph{4. Conditioning on the cardinality of the random set.}
For a uniformly random $\lfloor pn\rfloor$-element subset, the total
number of selected coordinates no longer fluctuates. The first projection
suggests centering the profile $h$ by its integral $1/2$, replacing $H$
by $H-1/4$ in one part of the calculation. This is only a heuristic:
a proof must also account for the global negative dependence introduced
by sampling without replacement.

\paragraph{5. Cyclic progressions and arithmetic obstructions.}
Determine the corrected avoidance law for progressions modulo $n$.
The interval degree profile is then replaced by a different incidence
geometry, and short cycles or small divisors of $n$ can affect overlap
counts. A result uniform in composite $n$ will require explicit handling
of these arithmetic obstructions.

\paragraph{6. Multicolor versions.}
For independent random colorings with prescribed color probabilities,
study the longest monochromatic progression and the joint head-count
vector by color. Same-site influences become vectors, and the relevant
coefficient should be a quadratic form rather than a single scalar.
Identify when cross-color exclusions contribute at leading order.

\paragraph{7. Waiting times and moments.}
Let $\tau_r=\min\{n:U_n\ge r\}$. Transfer the corrected distribution
to precise asymptotics for its quantiles, expectation, and variance.
Moment asymptotics require tail integration estimates beyond a local
inversion of Theorem~\ref{thm:critical}; the lattice effects must not be
averaged away without justification.

\paragraph{8. Kernel-checked formalization.}
Formalize the finite head bijection, same-difference incompatibility,
two-site rational-ratio lemma, and the one-coordinate covariance identity
first. Then formalize conditional local-lemma bounds and the finite
Bonferroni quotient expansion. The analytic degree-profile limit and the
lattice-maximization step can be separate modules. The code supplied here
is a regression suite, not a substitute for these proofs.

\section*{Conclusion}
The precise error of the longest-progression Poisson law is governed by
a geometric quantity that is invisible in an unweighted dependency-graph
bound: the squared entropy profile of progression incidence.
After separating the one-site projection from rare multiple overlaps,
the second factorial cumulant has an explicit positive coefficient.
A local avoidance expansion with a controlled connected-triple remainder
then turns this coefficient into a uniform probability correction.
The resulting order $(\log n)^2/n$ is sharp, and its leading constant
oscillates with a fully specified lattice phase. The next unresolved
scale in this analysis is the ratio-two contribution of order at most
$(\log n)^{3/2}/n$.

\clearpage
\bibliographystyle{plain}
\bibliography{references}
\end{document}
